\subsection{STRUCTURE OF CHAMBERS}

Let \(\mathbf{C}\) be a chamber, let \(\mathfrak{M}\) be the set of walls of \(\mathbf{C}\) , and for \(\mathrm{H} \in \mathfrak{M}\) let \(e_{\mathrm{H}}\) be the unit vector orthogonal to \(\mathrm{H}\) on the same side of \(\mathrm{H}\) as \(\mathbf{C}\) .

PROPOSITION 7. Assume that the group W is essential and finite. Then:

\begin{enumerate}
\def\labelenumi{(\roman{enumi})}
\item
  There exists a unique point \(a\) of \(\mathbf{E}\) invariant under \(\mathbf{W}\) .
\item
  The family \((e_{\mathsf{H}})_{\mathsf{H}\in \mathfrak{M}}\) is a basis of \(\mathrm{T}\) .
\item
  The chamber C is the open simplicial cone with vertex a defined by the basis \((e_{\mathrm{H}}^{\prime})_{\mathrm{H}\in \mathfrak{M}}\) of T such that \((e_{\mathrm{H}}|e_{\mathrm{H}'}') = \delta_{\mathrm{HH}'}\) .
\item
  By Prop. 4 of no. 6, there exists a point \(a \in E\) invariant under W. Let \(t \in T\) be such that \(t + a\) is invariant under W. For all \(w \in W\) ,
\end{enumerate}

\[
\mathrm{U} (w). t + a = w (t + a) = t + a,
\]

so \(\mathrm{U}(w).t = t\) ; since \(\mathrm{W}\) is essential, this implies that \(t = 0\) , showing the uniqueness of \(a\) .

\begin{enumerate}
\def\labelenumi{(\roman{enumi})}
\setcounter{enumi}{1}
\item
  Since W is essential, \(T = T_{1}\) in the notation of no. 7, and Prop. 5, (iv) shows that the family \((e_{\mathrm{H}})_{\mathrm{H}\in\mathfrak{M}}\) generates the vector space T. The existence of a point of E invariant under W shows that the family \((e_{\mathrm{H}})_{\mathrm{H}\in\mathfrak{M}}\) is free (no. 6, Prop. 4).
\item
  Let \(a\) be the unique point of \(\mathbf{E}\) invariant under \(\mathbf{W}\) . Since \((e_{\mathrm{H}})_{\mathrm{H}\in \mathfrak{M}}\) is a basis of \(\mathbf{T}\) , and since the scalar product is a non-degenerate bilinear form on \(\mathbf{T}\) , there exists a unique basis \((e_{\mathrm{H}}^{\prime})_{\mathrm{H}\in \mathfrak{M}}\) of \(\mathbf{T}\) such that \((e_{\mathrm{H}}|e_{\mathrm{H}'}') = \delta_{HH'}\) for \(\mathrm{H},\mathrm{H}'\) in \(\mathfrak{M}\) . Every point \(x\) of \(\mathbf{E}\) can be written uniquely in the form \(x = t + a\) with \(t = \sum_{\mathrm{H}\in \mathfrak{M}}\xi_{\mathrm{H}}.e_{\mathrm{H}}'\) and the \(\xi_{\mathrm{H}}\) real. Then \(x\) belongs to \(\mathbf{C}\) if and only if, for every hyperplane \(\mathrm{H}\in \mathfrak{M}\) , \(x\) is on the same side of \(\mathbf{H}\) as \(e_{\mathrm{H}}\) , or in other words \((t|e_{\mathrm{H}}) = \xi_{\mathrm{H}}\) is strictly positive. Hence (iii).
\end{enumerate}

PROPOSITION 8. Assume that the group W is essential, irreducible and infinite. Then:

\begin{enumerate}
\def\labelenumi{(\roman{enumi})}
\item
  No point of E is invariant under W.
\item
  We have \(\operatorname{Card} \mathfrak{M} = \dim T + 1\) , and there exist real numbers \(c_{\mathrm{H}} > 0\) such that \(\sum_{\mathrm{H} \in \mathfrak{M}} c_{\mathrm{H}} \cdot e_{\mathrm{H}} = 0\) . If the real numbers \(c_{\mathrm{H}}'\) are such that \(\sum_{\mathrm{H} \in \mathfrak{M}} c_{\mathrm{H}}' \cdot e_{\mathrm{H}} = 0\) , there exists a real number \(\xi\) such that \(c_{\mathrm{H}}' = \xi c_{\mathrm{H}}\) for all H in \(\mathfrak{M}\) .
\item
  The chamber C is an open simplex.
\end{enumerate}

Assertion (i) follows from Prop. 4. On the other hand, since W is essential, the vectors \((e_{\mathrm{H}})_{\mathrm{H}\in\mathfrak{M}}\) generate T. We have \((e_{\mathrm{H}}|e_{\mathrm{H}^{\prime}})\leqslant0\) for \(H,H^{\prime}\in \mathfrak{M}\) and \(H\neq H^{\prime}\) (Prop. 3) and, since W is irreducible, there does not exist a partition of \(\mathfrak{M}\) into two disjoint subsets \(\mathfrak{M}^{\prime}\) and \(\mathfrak{M}^{\prime\prime}\) such that \(H^{\prime}\in \mathfrak{M}^{\prime}\) and \(H^{\prime\prime}\in \mathfrak{M}^{\prime\prime}\) imply that \((e_{\mathrm{H}^{\prime}}|e_{\mathrm{H}^{\prime\prime}})=0\) . We can thus apply Lemma 5 of no. 5, and case 1) of that lemma is excluded; indeed, the \(e_{H}\) are not linearly independent, since W has no fixed point. Assertion (ii) follows.

We now prove (iii). Number the walls of C as \(H_{0}, H_{1}, \ldots, H_{d}\) and put \(t_{m} = e_{H_{m}}\) . By (ii), the vectors \(t_{1}, \ldots, t_{d}\) form a basis of T, so the hyperplanes \(H_{1}, \ldots, H_{d}\) have a point \(a_{0}\) in common, and there exists a basis \((t_{1}^{\prime}, \ldots, t_{d}^{\prime})\) of T such that \((t_{m}|t_{n}^{\prime}) = \delta_{mn}\) ; moreover, again by (ii) there exist real numbers \(c_{1} > 0, \ldots, c_{d} > 0\) such that

\[
t _ {0} = - \left(c _ {1}. t _ {1} + \dots + c _ {d}. t _ {d}\right).
\]

Since the vector \(t_{0}\) is orthogonal to the hyperplane \(H_{0}\) , there exists a real number c such that \(H_{0}\) is the set of points \(x = t + a_{0}\) of E with \((t|t_{0}) = -c\) .

Every point of E can be written uniquely in the form \(x = t + a_{0}\) with \(t = \xi_{1}.t_{1}^{\prime} + \cdots + \xi_{d}.t_{d}^{\prime}\) and \(\xi_{1}, \ldots, \xi_{d}\) real. The point x belongs to C if and only if it is on the same side of \(H_{m}\) as \(t_{m}\) for \(0 \leqslant m \leqslant d\) ; this translates into the inequalities \((t|t_{1}) > 0, \ldots, (t|t_{d}) > 0\) , and \((t|t_{0}) > -c\) , or equivalently by \(\xi_{1} > 0, \ldots, \xi_{d} > 0\) , \(c_{1}\xi_{1} + \cdots + c_{d}\xi_{d} < c\) . Since C is non-empty, c \textgreater{} 0. Put \(a_{m} = a_{0} + \frac{c}{c_{m}}.t_{m}'\) for \(1 \leqslant m \leqslant d\) ; then the chamber C consists of the points of E of the form \(a_{0} + \sum_{m=1}^{d} \lambda_{m}.(a_{m} - a_{0})\) with \(\lambda_{1} > 0, \ldots, \lambda_{d} > 0\) and \(\lambda_{1} + \cdots + \lambda_{d} < 1\) , so C is the open simplex with vertices \(a_{0}, \ldots, a_{d}\) . Q.E.D.

Remarks. 1) Identify E with \(E_{0} \times E_{1} \times \cdots \times E_{s}\) and W with \(W_{1} \times \cdots \times W_{s}\) as in no. 8. By Prop. 6, the chamber C is then identified with

\[
\mathrm{E} _ {0} \times \mathrm{C} _ {1} \times \dots \times \mathrm{C} _ {s},
\]

where \(C_{p}\) is a chamber in \(E_{p}\) with respect to the set of hyperplanes \(H_{p}\) . By Props. 7 and 8, each of the chambers \(C_{1},\ldots,C_{s}\) is an open simplicial cone or an open simplex.

\begin{enumerate}
\def\labelenumi{\arabic{enumi})}
\setcounter{enumi}{1}
\tightlist
\item
  Assume that W is irreducible and essential. If H and \(H'\) are two walls of C, then \(m_{HH'} = +\infty\) if and only if \(e_{H} = -e_{H'}\) (Prop. 3). By Props. 7 and 8, this can happen only if H and \(H'\) are the only walls of C and E is of dimension 1. Thus, the only case in which one of the \(m_{HH'}\) is infinite is that in which E is of dimension 1 and the group W is generated by the reflections associated to two distinct points (cf.~§2, no. 4).
\end{enumerate}

In the general case, the entries of the Coxeter matrix associated to W are finite unless at least one of \(E_{1},\ldots,E_{s}\) is of the preceding type.

